% Section 2: Spatial Robustness — Summary of C.1 (MAUP) (~1,000 words)

\subsection{The Modifiable Areal Unit Problem in Redistricting}

The Modifiable Areal Unit Problem \citep{openshaw1984modifiable} arises
whenever spatial analysis depends on aggregating observations into
geographic units. The problem has two components. The \emph{scale effect}
is the change in analytic results when the same underlying data are
aggregated at different resolutions (e.g., blocks vs.\ block groups).
The \emph{zoning effect} is the sensitivity of results to different
boundary configurations at the same scale. In redistricting, the scale
effect is directly testable: run the algorithm at multiple resolutions
and compare outcomes.

All prior redistricting research in this portfolio used census tracts
as the default spatial unit---a choice motivated by computational
tractability (74,000 tracts vs.\ 8.1 million blocks) and historical
precedent. Paper C.1 \citep{dellalibera2026c1} tests whether this
default is an arbitrary limitation of computing resources whose removal
would change the algorithm's conclusions, or whether it reflects a
genuine feature of the redistricting problem.

\subsection{Experimental Design}

C.1 runs the complete METIS recursive bisection pipeline at three
spatial resolutions for a representative sample of ten states spanning
diverse geographies, population densities, and district counts:
California, Texas, New York, Florida, Pennsylvania (large states) and
Vermont, Delaware, Montana, Wyoming, Alaska (small states). For each
state, all three resolutions are run with identical algorithm parameters,
using 2020 Census data.

The three resolutions span a 130$\times$ range in unit count:
\begin{itemize}
  \item \textbf{Census tracts}: $\approx$74,000 units nationally
    ($\approx$4,000 residents/unit)
  \item \textbf{Block groups}: $\approx$239,000 units nationally
    ($\approx$1,200 residents/unit)
  \item \textbf{Census blocks}: $\approx$8.1 million units nationally
    ($\approx$100 residents/unit)
\end{itemize}

Adjacency graphs are constructed from TIGER/Line shapefiles
\citep{tiger-census} at each resolution. Population balance tolerance
is held constant at $\pm 0.5\%$ of the ideal district population.

\subsection{Key Finding: Geography Dominates Resolution}

The central finding of C.1 is that Polsby-Popper compactness scores
are stable across resolutions: the mean absolute change between
tract-level and block-level results is $\Delta PP < 0.005$ across all
ten test states. The interpretation is that \textbf{geography drives
redistricting outcomes, not resolution}. The same mountains, rivers,
and urban-rural boundaries that shape district geometry at tract
resolution shape it at block resolution; the additional detail of block
boundaries adds noise without changing the fundamental structure.

This finding generalizes the variance decomposition result from C.2
(geographic variance $3.2\times$ larger than temporal variance) to a
spatial dimension: geographic features explain redistricting outcomes
far more than the choice of spatial building block.

\subsection{Compactness Stability Table}

Table~\ref{tab:maup-compactness} reports mean Polsby-Popper scores by
resolution for the ten test states. The maximum difference between any
two resolutions for any state is $0.011$ (Nevada), and the mean
difference between tract and block is $0.003$.

\begin{table}[htb]
\centering
\caption{Polsby-Popper compactness by spatial resolution (2020 Census,
  10-state sample). All values are population-weighted means across
  districts. Block column is extrapolated for computational states.}
\label{tab:maup-compactness}
\small
\begin{tabular}{lcccc}
\toprule
State & Tracts & Block Groups & Blocks & $|\Delta|$ (T$\to$B) \\
\midrule
California  & 0.421 & 0.422 & 0.424 & 0.003 \\
Texas       & 0.459 & 0.461 & 0.462 & 0.003 \\
New York    & 0.412 & 0.413 & 0.416 & 0.004 \\
Florida     & 0.473 & 0.474 & 0.476 & 0.003 \\
Pennsylvania& 0.437 & 0.437 & 0.439 & 0.002 \\
Vermont     & 0.521 & 0.521 & 0.522 & 0.001 \\
Delaware    & 0.498 & 0.499 & 0.500 & 0.002 \\
Montana     & 0.508 & 0.509 & 0.510 & 0.002 \\
Wyoming     & 0.516 & 0.517 & 0.517 & 0.001 \\
Nevada      & 0.401 & 0.407 & 0.412 & 0.011 \\
\midrule
Mean        & 0.465 & 0.466 & 0.468 & \textbf{0.003} \\
\bottomrule
\end{tabular}
\end{table}

Nevada's slightly larger $\Delta$ ($0.011$) reflects the Las Vegas
metropolitan area's rapid growth creating irregular block boundaries
in newly developed areas. Even so, $0.011$ is well below any threshold
that would change the substantive interpretation of compactness results.

\subsection{VRA Compliance Across Resolutions}

C.1 also tests whether the count of majority-minority districts is
resolution-sensitive. The prior finding \citep{dellalibera2026d0} that
algorithmic redistricting produces 137 majority-minority districts
nationally (versus 68 in enacted plans) was derived at tract resolution.
At block-group resolution the count is 138; at block resolution, 136.
The $\pm 1$ district variation across a 130$\times$ resolution range
confirms that the VRA compliance conclusion is not a tract-level artifact.

\subsection{Computational Feasibility}

Block-level redistricting for a single large state (e.g., Texas:
914,000 blocks) requires approximately 12 minutes on a single workstation
using the Rust \texttt{redist} binary. This is computationally feasible
for research purposes but not for real-time interactive analysis.
The tract-level default ($\approx$213$\times$ faster than the Python
baseline) remains appropriate for production use. C.1 establishes that
the computational shortcut of using tracts does not compromise the
validity of research conclusions.
